\section{Step 4: Solver settings for each speed}
\label{sec:solver}

Study~A uses the density-based solver for the two supersonic runs and the pressure-based solver at
$M = 0.95$, where the density-based solver did not converge.

\subsection{Density-based solver: supersonic runs}

\paragraph{How it works.} $\rho$, $\rho\mathbf u$ and $\rho E$ of every cell are updated together,
marching in pseudo-time until nothing changes; $k$ and $\omega$ are solved after them in each iteration.

\paragraph{Roe-FDS flux.} It splits the jump between two cells into pressure waves and flow, and takes each
from its upstream side: shocks come out sharp (2--3 cells) and without wiggles.

\paragraph{Courant number.}
\begin{equation}
\Delta t = \mathrm{CFL}\,\frac{\Delta x}{|u| + a}.
\end{equation}
CFL 1 while the shock forms, then 5; lower it to 3 if the residuals grow.

\paragraph{Order.} First-order upwind, $\phi_f = \phi_c$, is robust but smears the shock; second-order
upwind, $\phi_f = \phi_c + \nabla\phi\cdot\mathbf r$, uses least-squares cell-based gradients. Only
second-order results are used.

\paragraph{Steps.} \gui{General}: \gui{Density-Based}. \gui{Methods}: \gui{Formulation} = Implicit,
\gui{Flux Type} = Roe-FDS, \gui{Gradient} = Least Squares Cell Based, \gui{Flow}, \gui{Turbulent Kinetic
Energy} and \gui{Specific Dissipation Rate} = First Order Upwind (Second Order Upwind after the switch).
\gui{Controls}: \gui{Courant Number} = 1, then 5.

\paragraph{Why.} Shocks dominate and information only travels downstream, which is what the density-based
solver is built for. It converged in 3000--3200 iterations (Section~\ref{sec:runs}).

\subsection{Pressure-based solver: transonic run}

\paragraph{How it works.} The unknowns are velocity and pressure; the pressure equation comes from the mass
equation, and $\rho = p/(RT)$. \gui{Coupled} solves momentum and pressure together; the \gui{Pseudo Time
Method} damps the changes between iterations, with an automatic step.

\paragraph{Why at $M = 0.95$.} The flow is subsonic everywhere, waves travel upstream too, and the weak
shock at the shoulder moves with the pressure behind it. The density-based $C_D$ never settled:
\begin{itemize}
\item second order, CFL 3: 0.009--0.053;
\item non-reflecting outlet: 0.030--0.040;
\item CFL 1: 0.016--0.031.
\end{itemize}
The pressure-based solver converged in 1200 iterations, with residuals still falling ($10^{-6}$ to
$10^{-12}$).

\paragraph{Viscous Heating on.} Friction heats the boundary layer through the energy term
\begin{equation}
\Phi = \boldsymbol\tau : \nabla\mathbf u.
\end{equation}
The density-based solver always includes it; in the pressure-based solver it has to be ticked. Without it
the adiabatic wall stays near $T_\infty$ instead of $\approx T_0$, and the friction is wrong.

\paragraph{Residual criteria off.} Pressure-based residuals are scaled by their first iterations: they fell
below $10^{-3}$ at iteration 7 and Fluent stopped the run. Set all criteria to $10^{-12}$.

\paragraph{Steps.} \gui{General}: \gui{Pressure-Based}. \gui{Models}\then\gui{Viscous}\then
\gui{Options}: \gui{Viscous Heating}. \gui{Methods}: \gui{Scheme} = Coupled, \gui{Pseudo Time Method}
on (automatic time step), \gui{Gradient} = Least Squares Cell Based, \gui{Pressure} = Second Order,
\gui{Density}, \gui{Momentum}, \gui{Energy}, $k$ and $\omega$ = First Order Upwind (Second Order Upwind
after the switch). \gui{Monitors}\then\gui{Residuals}: all criteria $10^{-12}$.

\subsection{Summary and convergence criterion}

\begin{table}[H]
\centering
\caption{Solver settings for each speed.}
\label{tab:solvers}
\small
\begin{tabularx}{\linewidth}{@{}>{\raggedright\arraybackslash}p{2.3cm}>{\raggedright\arraybackslash}X>{\raggedright\arraybackslash}X@{}}
\toprule
 & Supersonic, $M = 1.2$ and 1.47 & Transonic, $M = 0.95$ \\
\midrule
Solver & density-based, implicit & pressure-based, Coupled, pseudo time \\
Main unknowns & density, momentum, energy, together & velocity and pressure, together; density from $p = \rho RT$ \\
Faces & Roe-FDS flux & upwind for density, momentum, energy, $k$, $\omega$; second-order pressure \\
Order & first order, then second order upwind & first order for 200 iterations, then second order \\
Gradients & least squares cell based & least squares cell based \\
Time step & CFL 1, then 5 (3 at the end of $M = 1.2$) & automatic pseudo time step \\
Viscous heating & always included & ticked by hand \\
Why & shock-dominated flow, information travels downstream & subsonic domain with a weak shock; the density-based solver swung \\
\bottomrule
\end{tabularx}
\end{table}

\paragraph{Convergence.}
\begin{equation}
\Delta = \frac{\max C_D - \min C_D}{\text{mean } C_D}\ \text{over the last 500 iterations} < 0.1\,\%.
\end{equation}
Residuals do not show whether the drag has settled.
